\chapter{Ambiente sperimentale}
\label{cha:ambienteSperimentale}

\section{Obiettivo e criteri di progettazione del testbed}

L'Obiettivo di base per la progettazione del testbed è quello di creare un ambiente
controllato, riproducibile e parametrizzabile all'interno del quale poter
valutare il differente comportamento dei demoni di sincronizzazione temporale in
condizioni di rete variabili. Nel caso in questione: linuxptp, chrony e ntpsec.
L'impostazione è stata scelta per mettere in evidenza non la precisione assoluta
dei nodi rispetto a una sorgente temporale esterna, quindi con un riferimento UTC,
ma la loro capacità nel mantenere una relazione temporale coerente all'interno
dell'ambiente sperimentale, al differire delle condizioni dei canali di comunicazione.

L'ambiente si basa su tre criteri fondamentali. Per primo bisogna isolare
completamente il testbed dalla rete esterna. Non essendoci necessità di riferimenti
esterni è stato possibile disabilitare ogni forma di connettività, garantendo il
massimo isolamento e costringendo i nodi a sincronizzarsi all'interno della sola
topologia di rete. Il secondo criterio è stato la riproducibilità dell'ambiente.
Infine l'ultimo ciriterio coinvolgeva la parametrizzazione degli impairment di
rete, i quali devono essere facilmetne modificabili, potendo gestire in modo
indipendente da ogni script la loro applicazione.

La topolgia è stata creata con l'utilizzo dei container Docker. I rami logici sono
stati opportunamenti separati e lo stesso ragionamento è stato applicato agli script,
con l'obiettivo di ridurre al minimo interferenze tra demoni e protocolli
differenti. Le condizioni di rete sono state modificate mediante l'utilizzo di Linux
Traffic Control.

\section{Topologia di rete}
La topologia si divide su due reti logiche sperate, rappresentative dell'evoluzione
della topologia T1, utilizzata per la fase di test preliminari, ma poi risultata
obsoleta. La rete A è stata associata a una subnet \texttt{10.0.10.0/24}, mentre la
rete B è stata associata ad una subnet \texttt{10.0.20.0/24}. I nodi \texttt{servergm},
\texttt{serverntp}, \texttt{clientchrony} e l'interfaccia \texttt{eth0} di \texttt{boundary}
sono stati collocati nella rete A, mentre \texttt{clientptp}, \texttt{clientntp}
e l'interfaccia \texttt{eth1} di\texttt{boundary} sono stati collocati nella rete
B. La scelta si basa sulla praticità operativa, si è cercato di evitare i conflitti
di configurazioni tra demoni e protocolli differenti, raccogliendo al meglio, sucessivamente,
i dati.

Il nodo \texttt{boundary} è stato configurato come nodo intermedio tra le due reti
assumendo un differente ruolo al variare del demone considerato. Nel caso di PTP,
assume un ruolo di boundary clock ricevendo sincronizzazione dalla rete A e propagandola
verso la rete B, mentrenel caso di NTPsec, è utilizzato come nodo intermedio tra \texttt{serverntp}
e \texttt{clientntp}. Per Chrony non è stato necessario il suo utilizzo.

{ \setlength{\intextsep}{8pt} \setlength{\textfloatsep}{8pt} \setlength{\floatsep}{8pt}

\begin{table}[H]\captionsetup{skip=4pt} \begin{tabular}{lll}\hline \textbf{Nodo} & \textbf{Rete} & \textbf{Ruolo sperimentale} \\ \hline \texttt{servergm} & A & Sorgente interna per PTP e Chrony \\ \texttt{serverntp} & A & Sorgente interna per NTPsec \\ \texttt{boundary} & A / B & Nodo intermedio; boundary clock PTP \\ \texttt{clientptp} & B & Nodo client PTP \\ \texttt{clientntp} & B & Nodo client NTPsec \\ \texttt{clientchrony} & A & Nodo client Chrony \\ \hline\end{tabular} \caption{Nodi della topologia T2 e ruoli assegnati nel testbed.} \label{tab:nodi_t2}\end{table}

\vspace{1cm}

\begin{figure}[H]\centering \begin{tikzpicture}[ node distance=1.5cm and 1.8cm, box/.style={ rectangle, draw, rounded corners, align=center, minimum width=3.8cm, minimum height=1.25cm, font=\small }, netbox/.style={ rectangle, draw, rounded corners, inner sep=0.45cm, fill=gray!8 }, arrow/.style={ -{Latex[length=3mm]}, thick }, edgeLabel/.style={ midway, fill=white, inner sep=1pt, font=\footnotesize } ]
% --- Nodi rete A ---
\node[box] (servergm) {%
\texttt{servergm}\\ Grandmaster PTP / Server Chrony\\ \texttt{eth0: 10.0.10.1} };

\node[box, right=2.2cm of servergm] (serverntp) {%
\texttt{serverntp}\\ Server NTPsec\\ \texttt{eth0: 10.0.10.4} };

\node[box, below=1.5cm of servergm] (clientchrony) {%
\texttt{clientchrony}\\ Client Chrony\\ \texttt{eth0: 10.0.10.3} };

\node[box, below=1.5cm of serverntp] (boundary) {%
\texttt{boundary}\\ Boundary Clock / nodo intermedio\\ \texttt{eth0: 10.0.10.2}\\ \texttt{eth1: 10.0.20.1} };


% --- Nodi rete B ---
\node[box, below=3.0cm of boundary, xshift=-2.2cm] (clientntp) {%
\texttt{clientntp}\\ Client NTPsec\\ \texttt{eth0: 10.0.20.2} };

\node[box, below=3.0cm of boundary, xshift=2.2cm] (clientptp) {%
\texttt{clientptp}\\ Client PTP\\ \texttt{eth0: 10.0.20.3} };


% --- Domini/reti ---
\begin{scope}[on background layer]\node[netbox, fit=(servergm)(serverntp)(clientchrony)(boundary)] (domainA) {};

\node[netbox, fit=(clientntp)(clientptp)] (domainB) {};\end{scope}


% --- Titoli domini ---
\node[font=\bfseries, above=0.15cm of domainA.north] {Rete A -- \texttt{10.0.10.0/24}};

\node[font=\bfseries, above=0.15cm of domainB.north] {Rete B -- \texttt{10.0.20.0/24}};


% --- Collegamenti logici ---
\draw[arrow] (servergm) -- node[edgeLabel] {Chrony} (clientchrony); \draw[arrow] (servergm) -- node[edgeLabel] {PTP} (boundary); \draw[arrow] (serverntp) -- node[edgeLabel] {NTPsec} (boundary); \draw[arrow] (boundary) -- node[edgeLabel, left] {NTPsec} (clientntp); \draw[arrow] (boundary) -- node[edgeLabel, right] {PTP} (clientptp);\end{tikzpicture} \caption{Topologia logica T2 utilizzata nel testbed sperimentale. La rete A ospita le sorgenti temporali interne e il client Chrony, mentre la rete B ospita i client PTP e NTPsec. Il nodo \texttt{boundary} collega le due reti tramite le interfacce \texttt{eth0} e \texttt{eth1}.} \label{fig:topologia_t2}\end{figure}

\vspace{0.4cm} }

Per ragioni configurative, si è reso necessario spearare la sorgente temporale
in due nodi differenti, distinguendo tra PTP e Chrony da un lato, e NTPsec dall'altro.
Conflitti di configurazione, installazione e permessi rendevano incompatibile l'installazione
di implementazioni software differenti sullo stesso nodo, soprattutto nei casi in
cui più servizi potevano tentare di intervenire sullo stesso clock locale.Inoltre
la separazione logica delle rete permette di avere una distinzione più chiara tra
le diverse catene di sincronizzazione analizzate nel testbed.

\section{Ambiente containerizzato}

Il testbed è realizzato tramite la creazione di una topologia di rete
containerizzata attraverso l'ausilio di Kathará, cioè un sisteama di emulazione di
rete open source basato su container. Questa scelta ha reso possibile la
riproduzione di una rete composta da un insieme di nodi Linux isolati, ognuno dotato
di una sua interfaccia, di file di configurazioni personalizzati e degli
specifici demoni di sincronizzazione. L'utilizzo dei container ha permesso inoltre
di automatizzare l'intero ciclo sperimentale, rendendo possibile l'inizializzazione
della topologia, la configuazione degli indirizzi IP, l'installazione dei
pacchetti necessari, la copia dei file di configurazione, l'avvio dei demoni, l'applicazione
degli impairment di rete e infine la raccolta dei log.

Ogni run può essere eseguita a partire da uno stato pulito della topologia,
riducendo il rischio che processi residui, configurazioni di rete precedenti o discipline
di coda già applicate influenzino le misurazioni successive.

\subsection{Kathará e organizzazione dei container}

La topologia T2 viene descritta nel file \texttt{lab.conf}, dove vengono
definiti i nodi del testbed, le differenti reti virtuali a cui essi appartengono,
i volumi condivisi montati all'interno dei container e le capability necessarie
all'esecuzione degli esperimenti.

Nella topologia, descritta come T2, i nodi sono associati a due reti logiche
distinte, indicate come rete A e rete B. I nodi \texttt{servergm}, \texttt{serverntp},
\texttt{clientchrony} e l'interfaccia \texttt{eth0} di \texttt{boundary} appartengono
alla rete A. I nodi \texttt{clientntp}, \texttt{clientptp} e l'interfaccia \texttt{eth1}
di \texttt{boundary} appartengono alla rete B. Il nodo \texttt{boundary} è l'unico
container connesso a entrambe le reti interne e costituisce il punto di interconnessione
tra i due domini.

\begin{lstlisting}[language=bash,caption={Estratto del file \texttt{lab.conf}: assegnazione dei nodi alle reti A e B.},label={lst:labconf_topologia}]

# Topologia T2

servergm[0]=A
serverntp[0]=A
boundary[0]=A
boundary[1]=B
clientchrony[0]=A
clientntp[0]=B
clientptp[0]=B
\end{lstlisting}

Nella fase iniziale di avvio, ogni contanier va ad eseguire il proprio file di inizializzazione
\texttt{.startup}. Essi contengono gli indirizzi IP delle interfacce, attivano i
link, popolano il file \texttt{/etc/hosts}, installano i pacchetti necessari e copiano
i file di configurazione dei demoni dalle directory condivise del progetto verso i
percorsi attesi dal sistema.

Nel nodo \texttt{boundary}, per esempio, si configurano due interfacce interne: \texttt{eth0}
sulla rete A e \texttt{eth1} sulla rete B.

\begin{lstlisting}[language=bash,caption={Estratto di \texttt{boundary.startup}: configurazione delle interfacce interne.},label={lst:boundary_startup_ip}]
ip addr add 10.0.10.2/24 dev eth0
ip addr add 10.0.20.1/24 dev eth1
ip link set eth0 up
ip link set eth1 up
\end{lstlisting}

Per quanto riguarda la risoluzione locale dei nomi, ogni nodo mantiene un file \texttt{/etc/hosts}
che associa gli hostname ai rispettivi indirizzi IP, evitando problematiche derivate
dalla dipendenza da servizi DNS esterni e rende esplicita la corrispondenza tra hostname
e indirizzi IP della topologia.

\begin{lstlisting}[language=bash,caption={Mappatura locale degli hostname nei container.},label={lst:hosts_mapping}]
10.0.10.1 servergm
10.0.10.4 serverntp
10.0.10.2 boundary0
10.0.20.1 boundary1
10.0.10.3 clientchrony
10.0.20.2 clientntp
10.0.20.3 clientptp
\end{lstlisting}

Distinguendo tra \texttt{boundary0} e \texttt{boundary1} si riesce a
identificare in modo chiaro le due interfacce del nodo \texttt{boundary}.
\texttt{boundary0} corrisponde all'interfaccia sulla rete A (\texttt{10.0.10.2}),
mentre \texttt{boundary1} corrisponde all'interfaccia sulla rete B (\texttt{10.0.20.1}).
Ciò è utile ai fini del testNTPsec, nei quali il nodo \texttt{clientntp}
utilizza il boundary come sorgente temporale raggiungibile sulla rete B.

All'interno del file \texttt{lab.conf} viene montata una directory del progetto all'interno
di tutti i container. Tale volume condiviso viene utilizzato sia per rendere
accessibili i file di configurazione ai nodi, sia per salvare i log prodotti durante
le run sperimentali.

\begin{lstlisting}[language=bash,caption={Estratto del file \texttt{lab.conf}: volume condiviso tra host e container.},label={lst:labconf_volume}]
servergm[volume]=/home/gabrielemenestrina/tesi_sync_lab|/tesi_sync_lab|rw
serverntp[volume]=/home/gabrielemenestrina/tesi_sync_lab|/tesi_sync_lab|rw
boundary[volume]=/home/gabrielemenestrina/tesi_sync_lab|/tesi_sync_lab|rw
clientchrony[volume]=/home/gabrielemenestrina/tesi_sync_lab|/tesi_sync_lab|rw
clientntp[volume]=/home/gabrielemenestrina/tesi_sync_lab|/tesi_sync_lab|rw
clientptp[volume]=/home/gabrielemenestrina/tesi_sync_lab|/tesi_sync_lab|rw
\end{lstlisting}

L'obiettivo alla base di tale scelta è duplice. Se da un lato permette di mantenere
tutti i file di configurazione all'interno della repository del progetto e copirali
nei differenti container in fase di avvio, allo stesso modo consente ai demoni e
agli script di scrivere i log in una directory persistente accessibile dall'host,
evitando problematicità dovute alla chiusura della topologia.

Ad esempio, il nodo \texttt{servergm} installa Chrony e linuxptp, copia la
configurazione Chrony lato server e predispone il file di configurazione PTP da
utilizzare come grandmaster interno del dominio sperimentale.

\begin{lstlisting}[language=bash,caption={Estratto di \texttt{servergm.startup}: installazione e copia delle configurazioni.},label={lst:servergm_startup}]
apt update -y
apt install -y chrony linuxptp

cp /tesi_sync_lab/configs/chrony/server_chrony.conf /etc/chrony/chrony.conf

pkill chronyd 2>/dev/null
sleep 2
chronyd &

mkdir -p /etc/ptp
cp /tesi_sync_lab/configs/ptp/server_ptp.conf /etc/ptp/server_gm.conf
\end{lstlisting}

Analogamente, il nodo \texttt{clientchrony} installa Chrony, copia la
configurazione client e avvia manualmente il demone \texttt{chronyd}. La scelta
dell'avvio manuale è stata preferita poichè, all'interno dei container, la
gestione tramite \texttt{systemctl} non è necessaria e può risultare meno adatta rispetto
all'avvio diretto del processo.

\begin{lstlisting}[language=bash,caption={Estratto di \texttt{clientchrony.startup}: configurazione e avvio del client Chrony.},label={lst:clientchrony_startup}]
apt update -y
apt install -y chrony

cp /tesi_sync_lab/configs/chrony/client_chrony.conf /etc/chrony/chrony.conf

pkill chronyd 2>/dev/null
sleep 2
chronyd &
\end{lstlisting}

In NTPsec, il nodo \texttt{serverntp} installa il pacchetto \texttt{ntpsec},
copia la configurazione server e avvia \texttt{ntpd}. Una procedura simile viene
adottata per il nodo\texttt{clientntp} che però utilizza una confiurazione lato
client. Per quanto riguarda invece il nodo \texttt{boundary} si ha sia la
configurazione PTP sia la configurazione NTPsec, ma i demoni non vengono mantenuti
attivi simultaneamente in quanto potrebbero essere causa di conflitti interni,
come problematicità di permessi.

\begin{lstlisting}[language=bash,caption={Estratto di \texttt{serverntp.startup}: configurazione e avvio di NTPsec.},label={lst:serverntp_startup}]
apt update -y
apt install -y ntpsec

cp /tesi_sync_lab/configs/ntp/server_ntp.conf /etc/ntpsec/ntp.conf

pkill ntpd 2>/dev/null
sleep 2
/usr/sbin/ntpd -g -c /etc/ntpsec/ntp.conf &
\end{lstlisting}

Il nodo \texttt{clientptp} installa invece \texttt{linuxptp} e copia il file di configurazione
PTP lato client. L'avvio effettivo del demone \texttt{ptp4l} non avviene direttamente
nel file \texttt{.startup}, ma viene gestito successivamente dagli script di esecuzione,
in modo da controllare l'ordine di avvio del grandmaster, del boundary clock e del
client PTP.

\begin{lstlisting}[language=bash,caption={Estratto di \texttt{clientptp.startup}: predisposizione della configurazione PTP lato client.},label={lst:clientptp_startup}]
apt update -y
apt install -y linuxptp

mkdir -p /etc/ptp
cp /tesi_sync_lab/configs/ptp/client_ptp.conf /etc/ptp/client_ptp.conf
\end{lstlisting}

Uno script differente, indicato come \texttt{bootstrapT3\_V2.sh}, viene utilizzato
per coordinare l'intera esecuzione sperimentale. Esso definisce gli scenari da
eseguire, l'insieme delle run, la directory di output dei log, il percorso della topologia
Kathará e l'interfaccia su cui applicare gli impairment per i test che perturbano
il ramo B.

\begin{lstlisting}[language=bash,caption={Variabili principali dello script \texttt{bootstrapT3\_V2.sh}.},label={lst:bootstrap_vars}]
SCENARIOS=("low" "medium" "high")
SCENARIO_DIR="$(dirname "$0")/scenarios"
RAWLOG_ROOT="analysis/raw_logs/T3_multiplerun"
NETNS="boundary"
IFACE="eth1"   # boundary lato rete B
TOPOLOGY_DIR="topologies/T2"

RUNS=15

mkdir -p "$RAWLOG_ROOT"
\end{lstlisting}

Ogni scenario viene eseguito più volte. Nel caso mostrato, per ciascuno scenario
sono previste quindici run indipendenti. Permettendo di ridurre eventuali rumori statistici
e di ottenere una distribuzione più affidabile dei risultati.

\subsection{Privilegi e capability}

Ogni esecuzione dei demoni di sincronizzazione richiede specifici privilegi all'interno
dei container, e lo stesso accade per l'applicazione degli impairment di rete introdotti,
da qui nasce la necessità di avviare in modalità privilegiata la topologia e di
assegnare capability aggiuntive ai nodi, \texttt{--privileged}. Così facendo si sono
ridotti problemi leagti alle restrizioni imposte dal runtime dei container.

Le capacibility sono un meccanismo del kernel Linux che consente di concedere privilegi
specifici a processi o container senza dover concedere loro l'accesso completo come
utente root. Di seguito sono riportate le capability principali utilizzate nel
testbed: \texttt{SYS\_TIME}, \texttt{NET\_RAW} e \texttt{NET\_ADMIN}.

\begin{lstlisting}[language=bash,caption={Estratto del file \texttt{lab.conf}: capability assegnate ai container.},label={lst:labconf_capabilities}]
clientptp[cap_add]=SYS_TIME,NET_RAW,NET_ADMIN
boundary[cap_add]=SYS_TIME,NET_RAW,NET_ADMIN
servergm[cap_add]=SYS_TIME,NET_RAW,NET_ADMIN
serverntp[cap_add]=SYS_TIME,NET_RAW,NET_ADMIN
clientchrony[cap_add]=SYS_TIME,NET_RAW,NET_ADMIN
\end{lstlisting}

\texttt{SYS\_TIME} consente processi nel container di intervenire sull'orologio di
sistema, questo indipendentemente se poi avviene o meno. Si è reso necessario l'aggiunta
di essa poiché era origine di conflitti internii, in quanto essa è necessaria
per i demoni che disciplinano il clock locale, come \texttt{chronyd}, \texttt{ntpd}
e \texttt{ptp4l}. La capability \texttt{NET\_ADMIN} è necessaria per operazioni amministrative
sulle interfacce di rete, tra cui l'applicazione di qdisc tramite \texttt{tc}.
La capability \texttt{NET\_RAW} abilita invece funzionalità legate all'uso di socket
raw e operazioni di rete richieste da alcuni strumenti o implementazioni.

\begin{lstlisting}[language=bash,caption={Avvio privilegiato della topologia nello script di automazione.},label={lst:kathara_privileged}]
sudo -v
printf 'n\n' | sudo -n kathara lstart -d "$TOPOLOGY_DIR" --privileged
\end{lstlisting}

Nella fase di testbed la topologia viene "ripulita" e riavviata prima di ogni
scenario o run. È stato effettuato tramite \texttt{kathara lclean} elimnando i
container della run precedente, evitando che processi attivi, configurazioni
residue o discipline di coda rimangano applicate all'esecuzione successiva.

\begin{lstlisting}[language=bash,caption={Reset della topologia prima dell'avvio di una run.},label={lst:kathara_lclean}]
kathara lclean -d "$TOPOLOGY_DIR" || true
sudo -v
printf 'n\n' | sudo -n kathara lstart -d "$TOPOLOGY_DIR" --privileged
\end{lstlisting}

L'uso di \texttt{|| true} impedisce che un eventuale errore nella fase di
pulizia, ad esempio dovuto all'assenza di una topologia già attiva, interrompa
lo script.

Un'ulteriore precauzione riguarda l'esecuzione dei demoni sul nodo \texttt{boundary}.
Poiché tale nodo può essere coinvolto sia nei test PTP sia nei test NTPsec, gli script
evitano di mantenere attivi contemporaneamente demoni che potrebbero tentare di disciplinare
lo stesso clock locale. Nel caso PTP, prima dell'avvio di \texttt{ptp4l} sul
grandmaster viene terminato \texttt{chronyd}; nel caso NTPsec, \texttt{ntpd}
viene avviato soltanto quando necessario e con PTP non attivo.

\begin{lstlisting}[language=bash,caption={Esempio di avvio controllato dei demoni PTP.},label={lst:ptp_start_sequence}]
kathara exec -d "$TOPOLOGY_DIR" servergm -- bash -lc 
"pkill chronyd 2>/dev/null || true; 
sleep 1; 
ptp4l -f /etc/ptp/server_gm.conf -m -i eth0 -S 
&> /tesi_sync_lab/$RUN_DIR/ptp_server.log &"

sleep 3

kathara exec -d "$TOPOLOGY_DIR" boundary -- bash -lc 
"ptp4l -f /etc/ptp/bc.conf -m -i eth0 -i eth1 -S 
&> /tesi_sync_lab/$RUN_DIR/ptp_boundary.log &"

sleep 3

kathara exec -d "$TOPOLOGY_DIR" clientptp -- bash -lc 
"ptp4l -f /etc/ptp/client_ptp.conf -m -i eth0 -S -s 
&> /tesi_sync_lab/$RUN_DIR/ptp_client.log &"
\end{lstlisting}

La sequenza di avvio è rilevante perché permette di inizializzare prima il grandmaster,
poi il boundary clock e infine il client PTP. In questo modo, il dominio PTP
viene costruito secondo la gerarchia attesa, riducendo la presenza di transitori dovuti
all'assenza iniziale di una sorgente temporale.

\subsection{Isolamento dalla rete esterna}

Nella fase di avvio si è reso necessario mantenere la connettività verso l'esterno
per consentire l'installazione dei pacchetti necessari, quali \texttt{chrony}, \texttt{ntpsec}
e \texttt{linuxptp}, e l'aggiornametno delle repository. Questo spiega la presenza
di bridge temporanei verso l'esterno, che vengono successivamente disabilitati
per garantire l'isolamento del testbed.

\begin{lstlisting}[language=bash,caption={Estratto del file \texttt{lab.conf}: bridge temporaneo verso l'esterno.},label={lst:labconf_bridged}]
servergm[bridged]=true
serverntp[bridged]=true
boundary[bridged]=true
clientchrony[bridged]=true
clientntp[bridged]=true
clientptp[bridged]=true
\end{lstlisting}

Successivamente, terminata l'inzializzazione della T2, la connettività esterna è stata
disabilitata utilizzando lo script \texttt{endInternetConnection.sh}.

Lo script disabilita l'interfaccia esterna sui nodi dotati di una sola
interfaccia interna e disabilita l'interfaccia esterna dedicata sul nodo \texttt{boundary}.
Nei nodi collegati a una sola rete interna, l'interfaccia \texttt{eth0} è
utilizzata per la topologia sperimentale, mentre \texttt{eth1} rappresenta l'interfaccia
bridged verso l'esterno. Nel nodo \texttt{boundary}, invece, \texttt{eth0} ed \texttt{eth1}
sono già utilizzate per le reti A e B; l'interfaccia esterna risulta quindi \texttt{eth2}.

\begin{lstlisting}[language=bash,caption={Script \texttt{endInternetConnection.sh}: disabilitazione della connettività esterna.},label={lst:end_internet}]
#!/bin/bash
source ./scripts/envT2.sh

echo "[+] Disattivo eth1 in tutti i nodi (post-start)"
for node in servergm serverntp clientchrony clientntp clientptp; do
if kathara exec -d "$TOPOLOGY_DIR" "$node" 'ip link set eth1 down' >/dev/null 2>&1; then
echo "[+] eth1 disattivata in $node"
else
echo "[!] eth1 non presente in $node"
fi
done

if kathara exec -d "$TOPOLOGY_DIR" boundary 'ip link set eth2 down' >/dev/null 2>&1; then
echo "[+] eth2 disattivata in boundary"
else
echo "[!] eth2 non presente in boundary"
fi

echo "[DONE] Connessione Internet disattivata in tutti i nodi."
\end{lstlisting}

Questo script viene eseguito in ogni run dopo l'avvio della topologia.

\begin{lstlisting}[language=bash,caption={Punto del workflow in cui viene disabilitata la rete esterna.},label={lst:end_internet_workflow}]
kathara lclean -d "$TOPOLOGY_DIR" || true
sudo -v
printf 'n\n' | sudo -n kathara lstart -d "$TOPOLOGY_DIR" --privileged

sleep 25

./scripts/endInternetConnection.sh

sleep 5
\end{lstlisting}

Togliendo connettività esterna, garantiamo che disturbi non controllati possono essere
introdotti all'interno del testbed, e allo stesso modo garantiamo che i
riferimenti temporali siano interni rispetto ai nodi preventivamente stabiliti, quali
sono rispettivametne i nodi \texttt{servergm} e \texttt{serverntp}.

\section{Configurazione logica dei protocolli}
Come è stato precedentemente descritto, vi sono tre percorsi logici distinti. Le differenti
configurazioni applicate ai container nella fase di avvio, vengono
opportunamente selezionate in base all'ambiente di test selezionato.

Le tre catene logiche considerate sono:
\begin{center}
  \texttt{servergm} $\rightarrow$ \texttt{boundary} $\rightarrow$ \texttt{clientptp}
\end{center}
per PTP,
\begin{center}
  \texttt{servergm} $\rightarrow$ \texttt{clientchrony}
\end{center}
per Chrony, e
\begin{center}
  \texttt{serverntp} $\rightarrow$ \texttt{boundary} $\rightarrow$ \texttt{clientntp}
\end{center}
per NTPsec.

\subsection{Configurazione PTP}

Nella configurazione PTP, costrutia tramtie il software linuxptp, in particolare
tramite il demone \texttt{ptp4l}, vi è il nodo \texttt{servergm} che opera come grandmaster
interno del dominio PTP, il nodo \texttt{boundary} che rappresenta il boundary clock
del sistema tra la rete A e la rete B, ed infinte si ha il \texttt{clientptp} ovvero
il nodo slave finale sulla rete B.

Nel dominio PTP si è utilizzato il dominio 24, ed in tutte le configurazione si
è adottato il meccanismo di delay end-to-end e il timestamping software, coerente
con la strutta dei container stessi, assenti di timestamping hardware dedicato.

La configurazione del nodo \texttt{servergm} è presente nel file \texttt{server\_ptp.conf}.

\begin{lstlisting}[language=bash,caption={Configurazione PTP del nodo servergm.},label={lst:ptp-servergm}]
[global]
twoStepFlag 1
domainNumber 24
delay_mechanism E2E
time_stamping software
slaveOnly 0
priority1 128
priority2 128
[eth0]
\end{lstlisting}

Il parametro \texttt{slaveOnly 0} rende possibile assegnare al nodo il ruolo di
master. I valori \texttt{priority1} e \texttt{priority2}, entrambi pari a 128, così
opportunamente scelti per \texttt{servergm}, permettono al nodo di avere prelazione
nella scelta come di selezione dei master. Così facendo si ha che la gerarchia
PTP viene resa coerente con il disegno logico sperimentale nel quale il \texttt{servergm}
rappresenta il riferimento temporale principale del dominio.

Differentemente si ha che il nodo \texttt{boundary} è configurato attraverso un
interfaccia multipla. È connesso alla rete A tramite \texttt{eth0} e alla rete B
tramite \texttt{eth1}. La configurazione è contenuta nel file \texttt{boundary\_ptp.conf}.

\begin{lstlisting}[language=bash,caption={Configurazione PTP del nodo boundary.},label={lst:ptp-boundary}]
[global]
twoStepFlag 1
domainNumber 24
delay_mechanism E2E
time_stamping software
slaveOnly 0
priority1 200
priority2 200
announceReceiptTimeout 3
tx_timestamp_timeout 20
[eth0]
[eth1]
\end{lstlisting}

Allo stesso modo del servergm anceh il \texttt{boundary} ha una configurazione di
tipo \texttt{slaveOnly 0}, dovendo svolgere il ruolo di master sulla porta downstream
verso \texttt{clientptp}. Tuttavia, i valori \texttt{priority1} e \texttt{priority2}
sono pari a 200, riducendo la preferenza rispetto al \texttt{servergm}. Questa configurazione
permette al boundary clock di sincronizzarsi rispetto al grandmaster sulla rete A
e di propagare la sincronizzazione verso la rete B.

Il nodo \texttt{clientptp} è configurato invece come slave PTP. La configurazione
è contenuta nel file \texttt{client\_ptp.conf}.

\begin{lstlisting}[language=bash,caption={Configurazione PTP del nodo clientptp.},label={lst:ptp-client}]
[global]
twoStepFlag 1
domainNumber 24
delay_mechanism E2E
time_stamping software
summary_interval 1
logging_level 6
slaveOnly 1
announceReceiptTimeout 3
tx_timestamp_timeout 20
verbose 1
logMinDelayReqInterval 0
logAnnounceInterval 1
logSyncInterval 0
[eth0]
\end{lstlisting}

Impostantdo il parametro \texttt{slaveOnly 1} si impedisce al client di
candidarsi come master, vincolandolo alla sincronizzazione rispetto al nodo upstream
selezionato nel dominio PTP. Per quanto rigurda i parametri rivolti agli
intervalli dei messaggi e al livello di logging, sono stati opportunamente modificati
per ottenre un comportamente osservabile più coerente con le neccessità dell'
analisi dati.

\subsection{Configurazione Chrony}

Nella configurazione Chrony il nodo \texttt{servergm} svolge il ruolo di
sorgente temporale interna, mentre \texttt{clientchrony} opera come client
Chrony.

La configurazione del server Chrony è contenuta nel file \texttt{server\_chrony.conf}.

\begin{lstlisting}[language=bash,caption={Configurazione Chrony del nodo servergm.},label={lst:chrony-server}]
server 127.127.1.0
local stratum 8
manual
makestep 1 -1
driftfile /var/lib/chrony/chrony.drift
bindaddress 0.0.0.0
allow 10.0.10.0/24
log measurements statistics tracking refclocks
logchange 0.01
\end{lstlisting}

La specifica \texttt{local stratum 8} permette al nodo di comportarsi come sorgente
temporale interna indipendentemente dalla presenza di un riferimento esterno. La
direttiva \texttt{allow 10.0.10.0/24} autorizza i nodi della rete A a interrogare
il server Chrony. La specifica \texttt{bindaddress 0.0.0.0} abilita l'ascolto sulle
interfacce disponibili del container.

Il client Chrony è configurato per utilizzare come sorgente il nodo \texttt{servergm},
raggiungibile all'indirizzo \texttt{10.0.10.1}. La configurazione è presente nel file
\texttt{client\_chrony.conf}.

\begin{lstlisting}[language=bash,caption={Configurazione Chrony del nodo clientchrony.},label={lst:chrony-client}]
server 10.0.10.1 iburst minpoll 1 maxpoll 2
driftfile /var/lib/chrony/chrony.drift
logdir /analysis/raw_logs/T2/chrony_logs
log measurements statistics tracking refclocks
logchange 0.01
\end{lstlisting}

L'opzione di \texttt{iburst} peremette al client del dominio di acquisire più
rapidamente la sorgente nella fase iniziale. I parametri \texttt{minpoll 1} e
\texttt{maxpoll 2} limitano l'intervallo di polling, mantenendo una frequenza di
campionamento elevata rispetto alla durata delle run sperimentali. La direttiva \texttt{driftfile}
specifica il file utilizzato da Chrony per memorizzare la stima del drift del
clock locale.

\subsection{Configurazione NTPsec}

Nella configurazione NTPsec, realizzata tramite \texttt{ntpd}, si ha come
riferimento temporale il nodo \texttt{serverntp}, utilizzando il clock locale. Il
nodo \texttt{boundary} si sincronizza con \texttt{serverntp} sulla rete A e rende
disponibile il servizio NTPsec verso la rete B vero il nodo \texttt{clientntp}, collocato
sulla rete B.

La catena logica è:

\begin{center}
  \texttt{serverntp} $\rightarrow$ \texttt{boundary} $\rightarrow$ \texttt{clientntp}
\end{center}

La configurazione è presente nel file \texttt{server\_ntp.conf}.

\begin{lstlisting}[language=bash,caption={Configurazione NTPsec del nodo serverntp.},label={lst:ntpsec-server}]
server 127.127.1.0
fudge 127.127.1.0 stratum 8
local stratum 8

driftfile /var/lib/ntpsec/ntp.drift

restrict default kod nomodify notrap nopeer limited
restrict 127.0.0.1
restrict ::1
restrict 10.0.10.0 mask 255.255.255.0 nomodify notrap
\end{lstlisting}

La sorgente \texttt{127.127.1.0} indica l'utilizzo del clock locale. Lo stratum
è impostato a 8, ciò è coerente con la sorgente temporale interna. La direttiva
\texttt{driftfile} specifica il file in cui NTPsec memorizza la stima del drift del
clock. Le direttive \texttt{restrict} limitano le operazioni consentite ai client
e autorizzano la rete \texttt{10.0.10.0/24} a interrogare il server.

La configurazione è contenuta nel file \texttt{boundary\_ntp.conf}.

\begin{lstlisting}[language=bash,caption={Configurazione NTPsec del nodo boundary.},label={lst:ntpsec-boundary}]
server 10.0.10.4 iburst
driftfile /var/lib/ntpsec/ntp.drift

restrict default kod nomodify notrap nopeer limited
restrict 10.0.10.0 mask 255.255.255.0 nomodify notrap
restrict 10.0.20.0 mask 255.255.255.0
restrict 127.0.0.1
restrict ::1
\end{lstlisting}

La direttiva \texttt{server 10.0.10.4 iburst} indica che il boundary si sincronizza
con \texttt{serverntp} sulla rete A. La rete B è autorizzata tramite la direttiva
\texttt{restrict 10.0.20.0 mask 255.255.255.0}, permettendo a \texttt{clientntp}
di utilizzare il boundary come sorgente NTPsec.

Il client NTPsec utilizza il boundary sulla rete B, raggiungibile all'indirizzo \texttt{10.0.20.1}.
La configurazione è contenuta nel file \texttt{client\_ntp.conf}.

\begin{lstlisting}[language=bash,caption={Configurazione NTPsec del nodo clientntp.},label={lst:ntpsec-client}]
server 10.0.20.1 iburst
driftfile /var/lib/ntpsec/ntp.drift
restrict 127.0.0.1
restrict ::1
\end{lstlisting}

Differentemente dalla configurazione PTP, il nodo \texttt{boundary} opera come classico
nodo intermedio e non come boundary clock nel senso del protocollo IEEE 1588.

\section{Emulazione degli impairment di rete}
\subsection{Linux Traffic Control}
\subsection{Parametri di degrado}
\subsection{Applicazione degli scenari}

\section{Automazione e raccolta dei log}
\subsection{Script di esecuzione}
\subsection{Organizzazione dei dati raccolti}